\subsection{Dinic: flujo máximo y corte mínimo}\label{rec:dinic}

Grafo dirigido con capacidades \codigo{long long}. Devuelve el flujo máximo.
Después del flujo, \codigo{cut(s)} marca el lado S del corte mínimo.

\tablafunc{
\codigo{Dinic(n)} & crea un grafo con \codigo{n} nodos & $O(n)$ \\
\codigo{add(u, v, cap)} & agrega arista dirigida con capacidad & $O(1)$ \\
\codigo{maxflow(s, t)} & valor del flujo máximo & $O(V^2E)$ \\
\codigo{cut(s)} & lado S después de saturar & $O(V+E)$ \\
\codigo{cut_edges(s)} & aristas originales que cruzan el corte & $O(V+E)$ \\
}

\copiar{el bloque \texttt{COPIA DINIC}. Usa \codigo{cut} y
\codigo{cut_edges} después de \codigo{maxflow}.}
\lstinputlisting{codigo/dinic.cpp}

\ojo{Para capacidades unitarias en una red bipartita, Hopcroft-Karp es más
simple. Para vértices con capacidad, divide cada vértice en entrada y salida
con una arista de capacidad.}